\newcommand{\tipoRequerimento}{Taxa administrativa da CIP}
\newcommand{\usaIPEstimada}{1}

\section{DOS FATOS E DOS DIREITOS}

\begin{enumerate}[resume]
\item É fato notório que a Resolução Normativa nº 888/2020 (REN nº 888/2020), promoveu alterações na redação da Resolução Normativa nº 414/2010 (REN nº 414/2010), acrescentando, entre outros, os seguintes dispositivos:

\Citacao{
    Seção VI
    
    Da Arrecadação da Contribuição para o Custeio do Serviço de Iluminação Pública

    \vspace{0.5cm}

    Art. 26-C. A contribuição para o custeio do serviço de iluminação pública, instituída pela legislação do poder municipal ou distrital, deve ser cobrada pelas distribuidoras nas faturas de energia elétrica nas condições previstas nessa legislação e demais atos normativos desses poderes.

    §1º A arrecadação de que trata o caput deve ser realizada pela distribuidora de forma não onerosa ao poder público municipal ou distrital. 

    §2º É vedado à distribuidora a realização da compensação dos valores arrecadados da contribuição com os créditos devidos pelo poder público municipal ou distrital, salvo quando houver autorização expressa na legislação municipal ou distrital. 
    §3º O repasse dos valores da contribuição para o custeio do serviço de iluminação pública deverá ocorrer até o décimo dia útil do mês subsequente ao de arrecadação, salvo disposição diversa na legislação e demais atos normativos do poder municipal ou distrital. 

    §4º A não observância dos §§2º e 3º implica a cobrança de multa de 2\% (dois por cento), atualização monetária e juros de mora previstos no art. 126, salvo disposição diversa na legislação e demais atos normativos do poder municipal ou distrital, sem prejuízo das sanções cabíveis.

    (grifou-se)
}

\item Ademais, a referida REN nº 888/2020 trouxe a seguinte previsão legal:

\Citacao{
    Art. 9º Estabelecer as seguintes datas-limites para as distribuidoras de energia elétrica adequarem os seus procedimentos às alterações promovidas por esta Resolução:

    §1º Até 13 de outubro de 2020, a distribuidora deverá: 

    I - notificar os municípios e o Distrito Federal sobre as alterações promovidas por esta Resolução, ressaltando as disposições dos arts. 26-C, 26-D e do §2º deste artigo em relação a contribuição para o custeio do serviço de iluminação pública e que o atual acordo operativo será substituído pela norma técnica de que trata o art. 21-A; e 

    II - encaminhar aos municípios e ao Distrito Federal os contratos que substituirão os contratos de iluminação pública e as novas minutas ou aditivos aos convênios e outros instrumentos celebrados, com as adequações necessárias ao disposto nesta Resolução. 

    §2º Faculta-se às distribuidoras a manutenção da cobrança pela arrecadação da contribuição para o custeio do serviço público de iluminação pública, no percentual máximo de 1 (um) por cento ou no percentual ora praticado, o que for menor, até a data de homologação de sua próxima revisão tarifária periódica, devendo a partir desta data cessar tal cobrança. 

    §3º Enquanto for mantida a cobrança pela distribuidora, de que trata o §2º, deverá ser realizada a reversão parcial das receitas auferidas para a propiciar a modicidade das tarifas do serviço de energia elétrica, conforme Submódulo 2.7 dos Procedimentos de Regulação Tarifária - PRORET. 

    Art. 10. Esta Resolução entra em vigor no dia 3 de agosto de 2020.

    (grifou-se)
}

\item Ocorre que, por meio do Mandado de Segurança Coletivo nº 1052154-94.2020.4.01.3400, a Associação Brasileira de Distribuidores de Energia Elétrica – ABRADEE buscou a suspensão de dispositivos da Resolução Normativa nº 888/2020, especificamente com o objetivo de afastar, até o julgamento definitivo do writ, os efeitos das modificações promovidas pela referida Resolução sobre os contratos firmados entre as distribuidoras associadas e os Municípios anteriormente à sua publicação.
\item Não obstante a impetração do referido mandado de segurança, e à luz do texto normativo vigente da REN nº 888/2020, a partir de 13/10/2020 (data da entrada em vigor) a distribuidora passou a estar vinculada ao limite máximo de 1\% (um por cento) sobre o valor mensalmente arrecadado a título de Contribuição para o Custeio do Serviço de Iluminação Pública (CIP/COSIP), a título de taxa de administração, até a data da primeira revisão tarifária subsequente.
\item A partir da homologação dessa revisão tarifária, o valor devido pela distribuidora a título de taxa de administração pela arrecadação da COSIP deveria ser reduzido a zero, vale dizer, a distribuidora deveria cessar, em caráter definitivo, qualquer cobrança de taxa de administração relativa à arrecadação da COSIP.
\item No caso em análise, a revisão tarifária periódica subsequente à entrada em vigor da REN nº 888/2020 ocorreu em 29/04/2021, nos termos da Resolução Homologatória nº 2.861/2021, de modo que essa data constitui o marco final para eventual cobrança de taxa de administração, nos estritos limites do § 2º do art. 9º da REN nº 888/2020.
Assim, a cobrança da taxa de administração da COSIP deveria observar, rigorosamente, os seguintes parâmetros normativos:

    \begin{enumerate}
        \item No período de 13/10/2020 a 29/04/2021, admitir-se-ia, tão somente, a cobrança de taxa de administração limitada ao percentual máximo de 1\% (um por cento) sobre o valor efetivamente arrecadado a título de CIP/COSIP.
        \item A partir de 29/04/2021 (data da revisão tarifária imediatamente subsequente à entrada em vigor da REN nº 888/2020), restou vedada qualquer cobrança de taxa de administração da COSIP pela distribuidora.
        \item Todo e qualquer valor cobrado a título de taxa de administração pela arrecadação da COSIP, em desacordo com os limites temporais e percentuais acima estabelecidos, deve ser integralmente restituído ao Município, com os correspondentes acréscimos legais (multa de 2\%, atualização monetária pelo IPCA e juros de mora de 1\% ao mês).
        \item Na hipótese de a taxa de administração da arrecadação da COSIP ter sido cobrada por meio da própria fatura de energia elétrica, a exação caracteriza cobrança indevida, impondo-se a devolução em dobro dos valores, nos termos do § 4º do art. 323 da REN nº 1.000/2021.
        
        \Citacao{
            Art. 323. A distribuidora, no caso de faturar valores incorretos, não apresentar fatura ou faturar sem utilizar a leitura do sistema de medição nos casos em que não haja previsão nesta Resolução, sem prejuízo das penalidades cabíveis, deve observar os seguintes procedimentos:

            [...]

            II - faturamento a maior: devolver ao consumidor e demais usuários, até o segundo ciclo de faturamento posterior à constatação, as quantias recebidas indevidamente nos últimos 60 ciclos de faturamento imediatamente anteriores à constatação. (Suspensos os efeitos referentes ao prazo de 60 (sessenta) ciclos, conforme DSP ANEEL 2.006, de 08.07.2024)

            § 2º No caso do inciso II do caput, a distribuidora deve devolver de acordo com as seguintes disposições:

            I - a quantia recebida indevidamente deve ser devolvida em dobro, independentemente de dolo ou culpa da distribuidora, salvo hipótese do §3º;

            II - o valor do inciso I deste parágrafo deve ser atualizado pelo Índice Nacional de Preços ao Consumidor Amplo – IPCA; e

            III - devem ser calculados e acrescidos os juros de mora à razão de 1\% ao mês pro rata die sobre o valor atualizado obtido do inciso II deste parágrafo.

            [...]

            § 4º A devolução em dobro prevista no § 2º é aplicável a todos os valores que compõem o faturamento, inclusive tributos, compensações, bandeiras tarifárias e cobranças de qualquer natureza.
        }
    \end{enumerate}


\item Registre-se, ainda, a recente renovação do Contrato de Concessão do Serviço Público de Distribuição de Energia Elétrica nº 26/2000-ANEEL, formalizada por meio da assinatura do \termoCIP{} Termo Aditivo pela \nomeConcessionaria{}. Referido aditivo, editado à luz das disposições do Decreto nº 12.068/2024, contém, em sua Cláusula Primeira, Subcláusula Quinta, disposição específica acerca da arrecadação de tributos por meio da fatura de energia elétrica, nos seguintes termos:

\Citacao{
    CLÁUSULA PRIMEIRA – OBJETO 

    Constitui objeto deste TERMO ADITIVO: 

    [...] 

    Subcláusula Quinta – A DISTRIBUIDORA aceita que a exploração do serviço público de distribuição de energia elétrica, de que é titular, seja realizada como função de utilidade pública prioritária, comprometendo-se a somente exercer outras atividades empresariais nos termos e condições previstas na legislação e na regulação da ANEEL, observando-se que: 

    [...] 

    IV. a arrecadação de tributos na fatura de energia elétrica decorrente de obrigação constitucional ou legal não será considerada atividade empresarial ou fonte de receitas alternativas, complementares e acessórias. (grifou-se)

}

\item Consta, ainda, da Subcláusula Sexta da Cláusula Primeira, obrigação expressa de observância da regulação expedida pela Aneel:

\Citacao{
    Subcláusula Sexta – Quaisquer normas, instruções, regulação ou determinações de caráter geral aplicáveis às prestadoras de serviço público de distribuição de energia elétrica, quando expedidas pelo PODERCONCEDENTE ou pela ANEEL, aplicar-se-ão automaticamente ao objeto da concessão ora contratada, a elas submetendo-se a DISTRIBUIDORA como condições implícitas e integrantes deste CONTRATO, observado o disposto no inciso V da Subcláusula Segunda da Cláusula Décima Quinta.
}

\item Ademais, a Cláusula Décima Oitava do referido Termo Aditivo estabelece que a distribuidora, além de aceitar as condições de prorrogação da concessão, assume, de forma expressa e irrevogável, a obrigação de desistir de eventuais ações judiciais individuais e de se abster de promover a execução ou de se valer dos efeitos de decisões proferidas em ações ajuizadas por entidades associativas, em representação, quando tais demandas contrariem as disposições do Termo Aditivo ou as condições fixadas pelo Decreto nº 12.068/2024, nos seguintes termos:

\Citacao{
    CLÁUSULA DÉCIMA OITAVA – DEMAIS DISPOSIÇÕES 

    [...] 
    
    Subcláusula Primeira - A DISTRIBUIDORA aceita na assinatura deste TERMO ADITIVO as condições de prorrogação estabelecidas no presente instrumento jurídico, bem como as disposições do Decreto nº12.068, de 2024. 
    
    Subcláusula Segunda - A DISTRIBUIDORA, de forma expressa e irrevogável, renuncia ao direito sobreo qual se fundam as ações, de qualquer natureza, que visem a questionar as cláusulas do presente TERMO ADITIVO bem como as condições do Decreto nº 12.068, de 2024. 
    
    Subcláusula Terceira - A DISTRIBUIDORA, de forma expressa e irrevogável, se compromete a desistir de eventuais ações judiciais individuais e a abster-se de promover a execução ou valer-se dos efeitos de decisões judiciais proferidas em ações ajuizadas por entidades associativas, em representação, cujo(s)direito(s) sobre os quais se funde a ação contrarie(m) as cláusulas do presente TERMO ADITIVO ou as condições do Decreto nº 12.068, de 2024.
}

\item Nesse contexto, verifica-se que se enquadram na hipótese da Subcláusula Terceira da Cláusula Décima Oitava do Termo Aditivo o Mandado de Segurança Coletivo nº 1052154-94.2020.4.01.3400, bem como as demais ações ajuizadas pela ABRADEE relacionadas ao atual art. 476 da REN nº 1.000/2021, que consolidou as disposições anteriormente veiculadas pela REN nº 888/2020.
\item Dessa forma, com a assinatura do Termo Aditivo e a consequente renovação do Contrato de Concessão, a \nomeConcessionaria{} está obrigada a:

    \begin{enumerate}
        \item Cumprir integralmente as disposições do art. 476 da REN nº 1.000/2021, assim como os artigos citados da REN 888/2020, cessando qualquer cobrança dirigida aos Municípios de sua área de concessão a título de remuneração pela arrecadação da CIP (ou COSIP) nas faturas de energia elétrica, bem como observar que eventual compensação prevista no § 2º do mesmo artigo somente poderá ser realizada se houver autorização expressa na legislação municipal pertinente; e
        \item Promover a desistência formal das ações judiciais, inclusive renunciando à representação processual nas demandas propostas pela ABRADEE que visem ao afastamento das disposições do art. 476 da REN nº 1.000/2021 (REN nº 888/2020), bem como efetivar as devoluções de valores devidos aos respectivos Municípios em decorrência dessa desistência.
    \end{enumerate}

\item À vista do exposto, e considerando a recente renovação do Contrato de Concessão e a assinatura do Termo Aditivo, impõe-se à distribuidora proceder à devolução dos valores cobrados indevidamente a maior, em decorrência da arrecadação da CIP em desconformidade com a Resolução Normativa nº 888/2020 e com a Resolução Normativa nº 1.000/2021, bem como em razão da obrigatória desistência das ações judiciais mencionadas, tudo em estrito cumprimento das obrigações contratuais assumidas como condição essencial e inafastável para a antecipação da renovação do contrato de concessão, nos termos do Decreto Federal nº 12.068/2024.

\end{enumerate}

\section{DOS PEDIDOS}

\begin{enumerate}[resume]

\item Diante de todo o exposto, requer-se à Concessionária que:

    \begin{enumerate}
        \item proceda à devolução dos valores indevidamente cobrados a maior a título de taxa de administração pela arrecadação da CIP, referentes ao período de 13/10/2020 até a data da efetiva e definitiva cessação da cobrança;
        \item O cálculo dos valores referidos no item anterior deverá observar, rigorosamente, os seguintes parâmetros: (i) no período compreendido entre 13/10/2020 e 29/04/2021, a cobrança da taxa de administração deverá ser limitada a 1\% (um por cento) do montante efetivamente arrecadado, sempre que o percentual contratualmente pactuado superar tal patamar; e (ii) a partir de 29/04/2021, o valor devido a esse título deve ser igual a zero, isto é, a cobrança da taxa de administração da CIP/COSIP deveria ter sido integralmente cessada naquela data.
        \item efetue a devolução dos valores mencionados na alínea anterior com os devidos acréscimos legais, compreendendo multa de 2\% (dois por cento), correção monetária pelo IPCA e juros de mora de 1\% (um por cento) ao mês, calculados pro rata die;
        \item realize a devolução, em dobro, dos valores apurados na forma da alínea anterior, haja vista que a cobrança da taxa de administração da CIP se deu por meio da fatura de energia elétrica, caracterizando cobrança indevida, nos termos do § 4º do art. 323 da REN nº 1.000/2021;
        \item proceda à restituição integral de todos os valores da CIP que tenham sido retidos para compensação de débitos do Município, com os correspondentes acréscimos legais acima descritos, uma vez que a legislação municipal não contém previsão expressa autorizando a distribuidora a realizar retenção de receitas da CIP para fins de compensação (encontro de contas) de dívidas municipais;
        \item assegure que a análise, a solução e a resposta à presente reclamação sejam proferidas dentro do prazo regulamentar previsto no art. 408, inciso I, da REN nº 1.000/2021, qual seja, 5 (cinco) dias úteis;
        \item promova a devolução dos valores apurados mediante depósito em conta corrente de titularidade do Município, nos termos do art. 323, § 7º, II, da REN nº 1.000/2021, com comunicação escrita do cumprimento das medidas dirigida ao representante que subscreve esta petição;
        \item determine que todas as comunicações, intimações e notificações relativas à presente reclamação sejam encaminhadas exclusivamente ao representante legal do Município que assina esta peça, preferencialmente para o endereço eletrônico \EmailEmpresa{}, em observância ao contraditório, à ampla defesa e aos princípios da eficiência e da celeridade processual, bem como ao art. 9º da REN nº 1.000/2021;

    \end{enumerate}

\end{enumerate}